\documentclass[11pt,twoside,twocolumn]{usgsreport}
\usepackage{usgsfonts}
\usepackage{usgsgeo}
\usepackage{usgsreporta}

\usepackage{graphicx}
\usepackage{enumitem}
\usepackage[T1]{fontenc}
\usepackage{etoolbox}

\usepackage{hyperref}
\hypersetup{
    pdftitle={GRIDGEN Release Notes},
    pdfauthor={U.S. Geological Survey},
    pdfsubject={GRIDGEN Release Notes},
    pdfkeywords={GRIDGEN, quadtree, unstructured grid, MODFLOW-USG, MODFLOW 6},
    bookmarksnumbered=true,
    bookmarksopen=true,
    bookmarksopenlevel=1,
    colorlinks=true,
    allcolors={blue},
    pdfstartview=Fit,
    pdfpagemode=UseOutlines,
    pdfpagelayout=TwoPageRight
}

\graphicspath{{./Figures/}}

\newcommand{\gridgenversion}{1.0.02}
\newcommand{\gridgendate}{January 6, 2017}

\renewcommand{\cooperator}
{George Mason University}
\renewcommand{\reporttitle}
{GRIDGEN: Release Notes}
\renewcommand{\coverphoto}{coverimage.jpg}
\renewcommand{\GSphotocredit}{Layered quadtree grid created with GRIDGEN.}
\renewcommand{\reportseries}{}
\renewcommand{\reportnumber}{}
\renewcommand{\reportyear}{2017}
\ifdef{\reportversion}{\renewcommand{\reportversion}{\gridgenversion}}{}
\renewcommand{\theauthors}{U.S. Geological Survey}
\renewcommand{\thetitlepageauthors}{Version \gridgenversion\ --- \gridgendate}
\renewcommand{\reportcitingtheauthors}{U.S. Geological Survey}
\renewcommand{\colophonmoreinfo}{}
\renewcommand{\reportbodypages}{}
\urlstyle{rm}
\renewcommand{\reportwebsiteroot}{https://doi.org/10.5066/}
\renewcommand{\reportwebsiteremainder}{F79G5JXV}
\ifdef{\usgsissn}{\renewcommand{\usgsissn}{}}{}
\renewcommand{\theconventions}{}
\definecolor{coverbar}{RGB}{0, 112, 60}
\renewcommand{\bannercolor}{\color{coverbar}}
\renewcommand{\reportrefname}{References Cited}

\begin{document}

\ifdef{\makefrontcoveralt}{\makefrontcoveralt}{\makefrontcover}
\ifdef{\makefrontmatterabv}{\makefrontmatterabv}{\makefrontmatter}

\onecolumn
\pagestyle{body}
\RaggedRight
\setlength{\parindent}{2em}

\section*{Abstract}
\addcontentsline{toc}{section}{Abstract}

GRIDGEN is a computer program for creating layered quadtree grids for use with numerical models, such as the MODFLOW-USG program for simulation of groundwater flow. The program begins by reading a three-dimensional base grid, which can have variable row and column widths and spatially variable cell top and bottom elevations. From this base grid, GRIDGEN will continuously divide into four any cell intersecting user-provided refinement features (points, lines, and polygons) until the desired level of refinement is reached. GRIDGEN will then smooth, or balance, the grid so that no two adjacent cells, including overlying and underlying cells, differ by more than a user-specified level tolerance. Once these gridding processes are completed, GRIDGEN saves a tree structure file so that the layered quadtree grid can be quickly reconstructed as needed. Once a tree structure file has been created, GRIDGEN can then be used to (1) export the layered quadtree grid as a shapefile, (2) export grid connectivity and cell information as ASCII text files for use with MODFLOW-USG or other numerical models, and (3) intersect the grid with shapefiles of points, lines, or polygons, and save intersection output as ASCII text files and shapefiles. The GRIDGEN program is demonstrated by creating a layered quadtree grid for the Biscayne aquifer in Miami-Dade County, Florida, using hydrologic features to control where refinement is added.

\section{Introduction}

This is GRIDGEN Version \gridgenversion. This GRIDGEN version is packaged for personal computers using the Microsoft Windows XP or 7 operating systems. Executable files for personal computers are provided as well as the source code. The source code can be compiled to run on other computers.

IMPORTANT: Users should review this document for a description of, and references for, this software. Users should also review the History section in this document, which describes changes that have been introduced into GRIDGEN with each official release; these changes may substantially affect users.

Version numbers for GRIDGEN will follow a major.minor.revision format. The major number will be increased when there are substantial new changes that break backward compatibility. The minor number will be increased when important, but relatively minor new functionality is added. The revision number will be added when errors are corrected in either the program or input files.

Instructions for installation, execution, and testing of GRIDGEN are provided below.

\section{Distribution File}

The following distribution file is for use on personal computers: gridgen1.0.02.zip. The distribution file contains:

\begin{itemize}
  \item Compiled Windows executables
  \item Documentation and text files
  \item Example problem
  \item Microsoft Visual Studio solution and project files
  \item Source code
\end{itemize}

The distribution file is a compressed zip file. The following directory structure is incorporated in the zip file:

\begin{verbatim}
|--gridgen.1.0.02
   |--bin       ; GRIDGEN executables for personal computers
   |--doc       ; Documentation files
   |--examples  ; Input files to run verification tests
   |--msvs      ; Microsoft Visual Studio solution and project files
   |--src       ; GRIDGEN source code for use on any computer
\end{verbatim}

It is recommended that no user files are kept in the gridgen directory structure. If you do plan to put your own files in the gridgen directory structure, do so only by creating additional subdirectories.

Included in the GRIDGEN distribution are various documentation files. Some of them are Portable Document Format (PDF) files. The PDF files are readable and printable on various computer platforms using Acrobat Reader from Adobe. The Acrobat Reader is freely available from the following World Wide Web site: \url{http://www.adobe.com/}

\section{Installation}

There is no installation of GRIDGEN other than the requirement that gridgen1.00.02.zip must be unzipped into a location where it can be accessed.

To make the executable versions of GRIDGEN accessible from any directory, the directory containing the executables should be included in the PATH environment variable. Also, if a prior release of GRIDGEN is installed on your system, the directory containing the executables for the prior release should be removed from the PATH environment variable.

As an alternative, the executable files, gridgen.exe and gridgen\_x64.exe, in the gridgen1.0.01/bin directory can be copied into a directory already included in the PATH environment variable.

\section{History}

This section describes changes introduced into GRIDGEN with each official release. These changes may substantially affect users.

\begin{itemize}
  % the release in progress, generated from items/*.toml
  \IfFileExists{develop.tex}{\input{develop.tex}}{}
  % release notes for GRIDGEN version 1.0.02
\item Version 1.0.02 01/6/2017
\begin{itemize}[label=--]
  \item The INTERPOLATE option for calculating top and bottom elevations was not working properly in some cases. This has been fixed.
  \item GRIDGEN would crash in some cases when saving a shared-vertex VTK file. This has been fixed.
  \item Polygons in the shapefile were not closed in the sense that the first vertex did not equal the last vertex. This has been fixed.
  \item A makefile has been added to the src folder for compiling on Linux and Mac operating systems.
  \item Added a ``Known Issues'' section to this document.
\end{itemize}

  % release notes for GRIDGEN version 1.0.01
\item Version 1.0.01 10/30/2015
\begin{itemize}[label=--]
  \item For models with very large x and y offsets, some of the numbers written to the output from GRIDGEN did not contain enough precision. The precision was increased for the values written to the *.nod and quadtree grid definition file.
  \item For grids that were rotated, the rotation origin was not being calculated correctly. This has been corrected. Note that the definition for ROTATION\_ANGLE in Table 2 of the GRIDGEN report is incorrect. The description for ROTATION\_ANGLE should be: ``Counter-clockwise rotation angle of lower left corner, in degrees.''
  \item One of the folders in the previous release did not zip correctly, and thus could not be extracted. This has been fixed.
\end{itemize}

  % release notes for GRIDGEN version 1.0.00
\item Version 1.0.00 2/6/2015
\begin{itemize}[label=--]
  \item This is the first public release of the GRIDGEN computer program. Details on this version of the program are in described in Lien and others (2014).
\end{itemize}

\end{itemize}

\section{System Requirements}

GRIDGEN is written in C++. The code has been used on UNIX-based computers and personal computers running various forms of the Microsoft Windows operating system.

\section{Executing GRIDGEN}

Two GRIDGEN executables for use on personal computers are provided. gridgen.exe is compiled for use on 32-bit Windows systems; gridgen\_x64 is compiled for use on 64-bit versions of Windows. In general, the gridgen\_x64 executable should be used in most instances.

GRIDGEN is initiated in a Windows Command-Prompt window by simply entering ``gridgen'' or ``gridgen\_x64'' followed by the appropriate command-line arguments, which are described in the user documentation.

\section{Testing}

The examples distributed with GRIDGEN can be run by navigating to the examples folder and executing the ``run.bat'' or ``run\_x64.bat'' batch files. These batch files will execute GRIDGEN multiple times in order to create a grid, create shape files, create input files for MODFLOW-USG, and intersect features with the resulting grid.

\section{Compiling}

The executable files provided in gridgen1.0.02/bin were created using the Microsoft Visual C++ compiler. Although executable versions of the program are provided, the source code is provided in the gridgen.1.0.02/src directory so that GRIDGEN can be recompiled if necessary. The Microsoft Visual Studio project files are also provided (in the msvs folder). However, the USGS cannot provide assistance to those compiling GRIDGEN. In general, the requirements are a C++ compiler and the knowledge of using the compiler.

\section{Test Problems}

The following is a list of test problems distributed with GRIDGEN.

\begin{itemize}
  \item Biscayne aquifer test problem -- This test problem is described in Lien and others (2014).
\end{itemize}

\section{Tips and Recommendations}

\begin{enumerate}
  \item Getting started---The best way to get started with GRIDGEN is to run the example problem that is included with the distribution. The example problem demonstrates the sequence of steps for generating a grid, intersecting the grid, and for writing grid shapefiles, MODFLOW-USG discretization input, and VTK files. The example files for GRIDGEN can then be used as a starting point for creating new applications.
  \item Vertical refinement considerations for MODFLOW-USG models---Testing with multi-layer models that have quad-based refinement has revealed oscillatory flow paths for some hydrogeologic conditions. These oscillatory flow paths are inconsistent with expected flow paths, and do not occur when the same horizontal grid is used for each layer. Use of a ghost-node correction in the vertical direction (which is not created by GRIDGEN Version 1.0.02; only horizontal ghost-node corrections are generated) can improve the simulated flows paths; however, the correction may not entirely remove the oscillations. For this reason, the current recommendation is to use the same horizontal grid for all of layers of a MODFLOW-USG model. This is achieved with GRIDGEN by listing the same set of refinement features for each model layer.
  \item For users of the PEST Groundwater Utilities, a grid specification file can be created from GRIDGEN output using the GRIDGEN2GSF program. GRIDGEN2GSF is included with the PEST Groundwater Utilities, available at: \url{http://www.pesthomepage.org/Groundwater_Utilities.php}
\end{enumerate}

\section{Known Issues}

\begin{itemize}
  \item When the base grid fits entirely within a specified active domain and the base grid does not touch the active domain, GRIDGEN excludes all cells. In this case, do not specify an active domain, and all cells will be included.
  \item GRIDGEN includes NODATA values in the surface interpolation when the ASCIIGRID option is used. Ensure that ASCIIGRID datasets do not have NODATA values in areas where surface interpolation is used.
  \item GRIDGEN writes a file with a ``.nod'' extension when a GRID\_TO\_USG block is executed. The first line in the .nod file has the number of cells and the number of connections. The number of connections does not include values for the node itself, and so this value does not directly correspond to the NJA value required by MODFLOW-USG. Instead, this value is equal to NJA - NODES.
\end{itemize}

\section{Documentation}

Lien, Jyh-Ming, Liu, Guilin, and Langevin, C.D., 2015, GRIDGEN version 1.0: A computer program for generating unstructured finite-volume grids: U.S. Geological Survey Open-File Report 2014-1109, 39 p., \url{http://dx.doi.org/10.3133/ofr20141109}.

\section{Disclaimer and Notices}

Please refer to the USGS Software User Rights Notice (\url{http://water.usgs.gov/software/help/notice/}) for complete use, copyright, and distribution information. The USGS provides no warranty, expressed or implied, as to the correctness of the furnished software or the suitability for any purpose. The software has been tested, but as with any complex software, there could be undetected errors. Users who find errors are requested to report them to the USGS.

References to non-USGS products, trade names, and (or) services are provided for information purposes only and do not constitute endorsement or warranty, express or implied, by the USGS, U.S. Department of Interior, or U.S. Government, as to their suitability, content, usefulness, functioning, completeness, or accuracy.

\section{Contacts}

\begin{verbatim}
U.S. Geological Survey
Office of Groundwater
411 National Center
Reston, VA 20192
(703) 648-5001
ogw_webmaster@usgs.gov
\end{verbatim}

\end{document}
